%%%%%%%%%%%%%%%%%%%%%%%%%%%%%%%%%%%%%%%%%
% Short Sectioned Assignment LaTeX Template Version 1.0 (5/5/12)
% This template has been downloaded from: http://www.LaTeXTemplates.com
% Original author:  Frits Wenneker (http://www.howtotex.com)
% License: CC BY-NC-SA 3.0 (http://creativecommons.org/licenses/by-nc-sa/3.0/)
%%%%%%%%%%%%%%%%%%%%%%%%%%%%%%%%%%%%%%%%%

% \documentclass[paper=a4, fontsize=11pt]{scrartcl}
\documentclass[11pt, a4paper]{book}
% Use 8-bit encoding that has 256 glyphs
\usepackage[T1]{fontenc}
\usepackage[utf8]{inputenc}
% Use the Adobe Utopia font for the document - comment this line to return to
% the LaTeX default
\usepackage{fourier}
% para insertar código con formato similar al editor
\usepackage{listings}
\usepackage[table,xcdraw]{xcolor}
\usepackage[most]{tcolorbox}
\usepackage[cache=false,outputdir=.]{minted}

\definecolor{grey}{RGB}{240,240,240}

%
% \begin{vcode}
%
% \end{vcode}
%
\newminted{py}{
    tabsize=10,
    linenos=true,
    breakanywhere,
    breaklines=true,
    mathescape,
    bgcolor=grey}


%
% \code{key to highlight}
%
\DeclareTotalTCBox{\code}{ O{grey} v O{} }{
   fontupper=\footnotesize\ttfamily,
   nobeforeafter,
   tcbox raise base,
   arc=0pt,
   top=3pt,
   bottom=3pt,
   left=3pt,
   right=3pt,
   leftrule=0pt,
   rightrule=0pt,
   toprule=0pt,
   bottomrule=0pt,
   boxsep=0.0mm,
   colback=#1,
   #3
}{#2}

% Dash ligatures
\usepackage{microtype}
\DisableLigatures[-]{}
% Selecciona el español para palabras introducidas automáticamente, p.ej.
% "septiembre" en la fecha y especifica que se use la palabra Tabla en vez de
% Cuadro
\usepackage[spanish, es-tabla]{babel}
% ,href} %para incluir URLs e hipervínculos dentro del texto (aunque hay que
% instalar href)
\usepackage{url}
% para incluir imágenes y colocarlas
\usepackage{graphics,graphicx, float}
\graphicspath{{../}{./}}
% para incluir el símbolo del euro
\usepackage[gen]{eurosym}
% para incluir citas del archivo <nombre>.bib
\usepackage{cite}
\usepackage{enumerate}
\usepackage{hyperref}
\usepackage{graphicx}
\usepackage{tabularx}
\usepackage{multirow}
\usepackage{booktabs}

\newcommand{\mc}[3]{\multicolumn{#1}{#2}{#3}}
\newcommand{\mr}[3]{\multirow{#1}{#2}{#3}}
\newcommand{\nl}{\tabularnewline}

\hypersetup{
    colorlinks=true,    % false: boxed links; true: colored links
    linkcolor=black,    % color of internal links
    urlcolor=cyan       % color of external links
}
\renewcommand{\familydefault}{\sfdefault}
% Custom headers and footers
\usepackage{fancyhdr}
% Makes all pages in the document conform to the custom headers and footers
\pagestyle{fancyplain}
% Empty left header
\fancyhead[L]{}
% Empty center header
\fancyhead[C]{}
% My name
\fancyhead[R]{Luis Liñán Villafranca}
% Empty left footer
\fancyfoot[L]{}
% Empty center footer
\fancyfoot[C]{}
% Page numbering for right footer
\fancyfoot[R]{\thepage}
% Remove header underlines
%\renewcommand{\headrulewidth}{0pt}
% Remove footer underlines
\renewcommand{\footrulewidth}{0pt}
% Customize the height of the header
\setlength{\headheight}{13.6pt}

\usepackage{titlesec, blindtext, color}
\definecolor{gray75}{gray}{0.75}
\newcommand{\hsp}{\hspace{20pt}}
\titleformat{\chapter}[hang]{%
    \Huge\bfseries%
}{\thechapter\hsp\textcolor{gray75}{|}\hsp}{0pt}{\Huge\bfseries}
\setcounter{secnumdepth}{4}
\usepackage[Lenny]{fncychap}
\usepackage{minted}
\usepackage{tikz}
\usetikzlibrary{babel,positioning}
\usepackage{calc}
\usepackage{caption}

\begin{document}

    % Plantilla portada UGR
    \begin{titlepage}
\newlength{\centeroffset}
\setlength{\centeroffset}{-0.5\oddsidemargin}
\addtolength{\centeroffset}{0.5\evensidemargin}
\thispagestyle{empty}

\noindent\hspace*{\centeroffset}\begin{minipage}{\textwidth}

\centering
\includegraphics[width=0.9\textwidth]{logos/logo_ugr.jpg}\\[1.4cm]

\textsc{ \Large TRABAJO FIN DE GRADO\\[0.2cm]}
\textsc{ GRADO EN INGENIERIA INFORMATICA}\\[1cm]

{\Huge\bfseries Título \\}
\noindent\rule[-1ex]{\textwidth}{3pt}\\[3.5ex]
{\large\bfseries Subtítulo }
\end{minipage}

\vspace{2.5cm}
\noindent\hspace*{\centeroffset}
\begin{minipage}{\textwidth}
\centering

\textbf{Autor}\\ {Luis Liñán Villafranca}\\[2.5ex]
\textbf{Director}\\ {Juan Julián Merelo Guervós}\\[2cm]
\includegraphics[width=0.3\textwidth]{logos/etsiit_logo.png}\\[0.1cm]
\textsc{Escuela Técnica Superior de Ingenierías Informática y de
Telecomunicación}\\
\textsc{---}\\
Granada, Septiembre de 2019
\end{minipage}
\end{titlepage}


    % Plantilla prefacio UGR
    \thispagestyle{empty}

\begin{center}

{\large\bfseries Deep G. Prop. \\ Optimización de parámetros de perceptrones
multicapa}\\

\end{center}

\begin{center}

    Luis Liñán Villafranca \\

\end{center}

\vspace{0.7cm}

\vspace{0.5cm}
\noindent{\textbf{Palabras clave}:
\textit{Software libre},
\textit{Redes neuronales},
\textit{Algoritmos genéticos},
\textit{Optimización de parámetros},
\textit{Clasificación}}
\vspace{0.7cm}

\noindent{\textbf{Resumen}}\\

Existe una cuestión que concierne a los modelos de clasificación supervisada:
encontrar la configuración de parámetros adecuada al problema al que se aplica.
Los perceptrones multicapa (MLP) con propagación hacia atrás (backpropagation)
tienen multitud valores de entrada, como el número de capas ocultas, el número
de neuronas (o perceptrones) en cada capa, los pesos entre las neuronas, la
función de activación de cada neurona... Lo cual hace que su correcta
inicialización suponga una tarea difícil y costosa. Los algoritmos genéticos
(GA) son una buena opción para solventar ésta, ya que realizando las
operaciones características de éstos (que simulan el proceso de selección
natural) se puede obtener una solución optimizada de parámetros que sirvan como
punto de partida para los perceptrones multicapa.

\cleardoublepage

\begin{center}

    {\large\bfseries Deep G. Prop. \\ Multilayer perceptrons parameters
    optimization}\\

\end{center}

\begin{center}

    Luis Liñán Villafranca\\

\end{center}

\vspace{0.5cm}
\noindent{\textbf{Keywords}:
\textit{free software}
\textit{neural networks},
\textit{genetic algorithms},
\textit{parameters optimization},
\textit{classification}}
\vspace{0.7cm}

\noindent{\textbf{Abstract}}\\

Supervised learning (categorized inside machine learning) deals with classified
data, in which each example comes associated to a label. The goal of these
problems is to train a model to be able to correctly classify a new example
without knowing the associated label. One of the big problems of supervised
classification is the ability to find the appropriate parameters configuration
for a particular dataset and model which enables the correct classification of
most of the examples. Multilayer perceptrons (MLP) have a huge variety of
initial parameters to configure the neural network with such as the number of
hidden layers, the number of neurons for each layer, the connection weights
between neurons, the activation function for each neuron... They also have
other parameters associated to weights training when using backpropagation as
learning method such as the optimizer or the loss function. The big quantity of
these aforementioned parameters makes the search a difficult and expensive
undertaking. Among the big variety of methods and algorithms present out there,
the Genetic algorithms (GA) are a great choice. Making use of its common
operations mimicking natural selection flow (selection, crossover, mutation and
replacement) to settle this task, may obtain optimized initial parameters for
the multilayer perceptrons future application.


\cleardoublepage

\thispagestyle{empty}

\noindent\rule[-1ex]{\textwidth}{2pt}\\[4.5ex]

D. \textbf{Juan Julián Merelo Guervós}, Profesor(a) del ...

\vspace{0.5cm}

\textbf{Informo:}

\vspace{0.5cm}

Que el presente trabajo, titulado \textit{\textbf{Deep G. Prop.}}, ha sido
realizado bajo mi supervisión por \textbf{Luis Liñán Villafranca}, y autorizo
la defensa de dicho trabajo ante el tribunal que corresponda.

\vspace{0.5cm}

Y para que conste, expiden y firman el presente informe en Granada a Junio de
2020.

\vspace{1cm}

\textbf{El/la director(a)/es:}

\vspace{5cm}

\noindent\textbf{Juan Julián Merelo Guervós}

\chapter*{Agradecimientos}






    % Índice de contenidos
    \newpage
    \tableofcontents

    % Índice de imágenes y tablas
    \newpage
    \listoffigures

    % Si hay suficientes se incluirá dicho índice
    \listoftables
    \newpage

    % Introducción
    \chapter{Introducción}



    % Descripción del problema y hasta donde se llega
    \chapter{Descripción del problema}



    % Estado del arte
    % 1. Crítica al estado del arte
    % 2. Propuesta
    \chapter{Estado del arte}



    % Planificación del trabajo
    \chapter{Planificación}

\section{Metodología utilizada}

\section{Temporización}

\section{Seguimiento del desarrollo}


    % Análisis del problema
    % 1. Análisis de requisitos
    % 2. Análisis de las soluciones
    % 3. Solucion propuesta
    % 4. Análisis de seguridad
    \chapter{Análisis del problema}
 


    % Desarrollo bajo sprints:
    % 1. Permitir registros y login de usuarios
    % 2. Desarrollo del sistema de incidencias
    % 3. Desarrollo del sistema de denuncias administrativas y accidentes
    % 4. Desarrollo del sistema de croquis
    % 5. Instalación de la aplicación de manera automática
    \chapter{Implementación}

La implementación del software se ha dividido en hitos. Estos, han sido
definidos en Github y cada uno de ellos contiene un grupo de \textit{issues} que
se corresponden con las distintas mejoras que se han ido incorporando al
software a lo largo de su desarrollo.\\



    % Presupuesto

    % Conclusiones
    \chapter{Conclusiones y trabajos futuros}



    % Trabajos futuros

    \newpage
    \bibliography{bibliografia}
    \bibliographystyle{plain}

\end{document}
